% Este é o modelo do abnTeX2
% ---
% inserir lista de abreviaturas e siglas
% ---
%\begin{siglas}
%  \item[ABNT] Associação Brasileira de Normas Técnicas
%  \item[abnTeX] ABsurdas Normas para TeX
%\end{siglas}
% ---

%
% Outra opção é utilizar o pacote Acronym (mais funcional).
%

%\ac{SW}  	SW documents are beautifully typeset.
%\acf{SW} 	Scientific Word (SW) documents are beautifully typeset.
%\acs{SW} 	SW documents are beautifully typeset.
%\acl{SW} 	Scientific Word documents are beautifully typeset.

% Este é o modelo do abnTeX2
% ---
% inserir lista de abreviaturas e siglas
% ---
%\begin{siglas}
%  \item[ABNT] Associação Brasileira de Normas Técnicas
%  \item[abnTeX] ABsurdas Normas para TeX
%\end{siglas}
% ---

%
% Outra opção é utilizar o pacote Acronym (mais funcional).
%

%\ac{SW}  	SW documents are beautifully typeset.
%\acf{SW} 	Scientific Word (SW) documents are beautifully typeset.
%\acs{SW} 	SW documents are beautifully typeset.
%\acl{SW} 	Scientific Word documents are beautifully typeset.

\pdfbookmark[0]{\listadesiglasname}{los}
\chapter*{\listadesiglasname}%

% Incluir no sumário
% \addcontentsline{toc}{chapter}{\listadesiglasname}%

\begin{acronym}
  \acro{API}{\textit{Application Programming Interface}}
  \acro{AP}{\textit{Access Point}}
  \acro{CCQ}{\textit{Client Connection Quality}}
  \acro{COSI}{Colegiado do Curso de Sistemas de Informação}
  \acro{DECSI}{Departamento de Computação e Sistemas}
  \acro{HTTP}{\textit{Hypertext Transfer Protocol}}
  \acro{ICEA}{Instituto de Ciências Exatas e Aplicadas}
  \acro{JM}{João Monlevade}
  \acro{RF}{Rádio Frequência}
  \acro{SI}{Sistemas de Informação}
  \acro{SNMP}{\textit{Simple Network Management Protocol}}
  \acro{UFOP}{Universidade Federal de Ouro Preto}
  \acro{VRP}{\textit{Vehicle Routing Problem}}
  \acro{WLAN}{\textit{Wireless Local Area Network}}
\end{acronym}
